% !TeX program = xelatex
\documentclass[11pt]{article}
\newcommand{\notesdir}{../../lecture-notes}
\providecommand{\notesdir}{.}

% -----------------------------------------------------------------------------
% Shared preamble for the CO5059 lecture notes and the quiz-review documents.
% Documents outside this folder define \notesdir (path to this folder) first.
% -----------------------------------------------------------------------------

% Page layout and typography
\usepackage[a4paper,portrait,margin=1in]{geometry}
\usepackage{fontspec}
\newfontfamily\arialfont{Arial}
\usepackage{microtype}
\usepackage{parskip}

% Mathematics and diagrams
\usepackage{amsmath,amssymb,amsthm,mathtools}
\usepackage{tikz-cd}
\usetikzlibrary{arrows.meta,positioning,shapes.geometric,shapes.arrows,calc,fit,backgrounds}

% Reusable styles for flowcharts and concept diagrams.
\tikzset{
  concept/.style={
    draw=brandprimary,
    fill=brandsecondary!8,
    rounded corners,
    minimum width=3.2cm,
    minimum height=0.9cm,
    align=center,
    font=\small,
  },
  diagram arrow/.style={-{Latex[length=2.5mm]},thick,brandprimary},
}

% Tables, figures, code, and lists
\usepackage{graphicx}
\graphicspath{{\notesdir/}}
\usepackage{enumitem}
\usepackage{booktabs}
\usepackage{array}
\usepackage{listings}

% Color and links
\usepackage[dvipsnames]{xcolor}
\definecolor{brandprimary}{HTML}{030391}
\definecolor{brandsecondary}{HTML}{1488D8}
\usepackage[
  colorlinks=true,
  linkcolor=brandprimary,
  citecolor=brandsecondary,
  urlcolor=brandsecondary
]{hyperref}

% Headers and footers
\usepackage{fancyhdr}
\setlength{\headheight}{36pt}
\pagestyle{fancy}
\fancyhf{}
\fancyhead[L]{%
  \parbox[c][0.42in][c]{0.42\textwidth}{%
    \raggedright\small\textbf{\coursecode}\\\coursename}}
\fancyhead[R]{%
  \parbox[c][0.42in][c]{0.42\textwidth}{%
    \raggedleft\includegraphics[height=0.42in]{assets/hcmut_logo.png}}}
\fancyfoot[L]{\small Semester \semester}
\fancyfoot[C]{\small\nouppercase{\leftmark}}
\fancyfoot[R]{\small Page \thepage}
\renewcommand{\sectionmark}[1]{\markboth{#1}{}}

% -----------------------------------------------------------------------------
% Course information -- change these values once per course
% -----------------------------------------------------------------------------
\newcommand{\universityname}{HCMC University of Technology}
\newcommand{\faculty}{Faculty of Computer Science and Engineering}
\newcommand{\coursecode}{CO5059}
\newcommand{\coursename}{Advanced System Architectures}
\newcommand{\semester}{261}
\newcommand{\instructor}{Tran Ngoc Thinh}
\newcommand{\instructoremail}{}% leave empty to hide
\newcommand{\notetaker}{Phan Nhat Minh}
\newcommand{\notetakeremail}{minh.phannhat95@hcmut.edu.vn}
\newcommand{\coursedescription}{Quantitative design and evaluation of computer
systems: performance and cost, instruction set architectures, pipelining and
its hazards, memory hierarchy, multiprocessors, superscalar/VLIW processors,
and AI \& security chips.}
\newcommand{\emaillink}[1]{%
  \href{mailto:#1}{\textcolor{brandprimary}{\texttt{#1}}}}

% Colored definitions, results, examples, proofs, and notes.
\input{\notesdir/colored-boxes}

% -----------------------------------------------------------------------------
% Course-specific notation
% -----------------------------------------------------------------------------
\newcommand{\CPI}{\mathrm{CPI}}

\lstdefinestyle{mips}{
  basicstyle=\ttfamily\small,
  keywordstyle=\bfseries\color{brandprimary},
  morekeywords={lw,sw,add,sub,and,or,slt,beq,bne,beqz,xor,nop,addi,addiu,addu,
    subu,li,la,jal,jr,jalr,sll,srl,sra,sllv,srlv,srav,andi,ori,xori,nor,slti,
    sltiu,sltu,lui,mult,multu,div,divu,mfhi,mflo,ll,sc,syscall,move,lh,lhu,lb,
    lbu,sh,sb,blez,bgtz,bltz,bgez,j,mul},
  showstringspaces=false,
  columns=fullflexible,
  keepspaces=true,
}
\lstset{style=mips}

% Creates a consistently formatted lecture heading.
% Usage: \lecture{number}{title}{date}
\newcommand{\lecture}[3]{%
  \section{Lecture #1: #2}
  \label{lec:#1}
  \begin{center}
    \small\textit{#3}
  \end{center}
}


% Cover page + contents page.  Usage: \makecover{Big title}{Description}
\newcommand{\makecover}[2]{%
\pagenumbering{roman}
\begin{titlepage}
  \thispagestyle{empty}

  \begin{center}
    \includegraphics[width=0.82\textwidth]{assets/hcmut_banner_en.png}

    \vspace{0.35cm}
    {\arialfont\large\bfseries\color{brandprimary}\faculty\par}
  \end{center}

  \vfill
  \begin{center}
    {\Huge\bfseries #1\par}
    \vspace{0.55cm}
    {\LARGE\bfseries\color{brandprimary}\coursename\par}
    \vspace{0.2cm}
    {\arialfont\large\bfseries\color{brandsecondary}\coursecode\par}

    \vspace{0.65cm}
    {\color{brandsecondary}\rule{0.38\textwidth}{1.2pt}}
    \vspace{0.65cm}

    {\large\textbf{Instructor:} \instructor
      \ifx\instructoremail\empty\else\ {\small(\emaillink{\instructoremail})}\fi\par}
    \vspace{0.3cm}
    {\large\textbf{Notes by:} \notetaker\
      {\small(\emaillink{\notetakeremail})}\par}
    \vspace{0.3cm}
    {\large\textbf{Semester:} \semester\par}

    \vspace{0.8cm}
    \begin{minipage}{0.72\textwidth}
      \centering\itshape\color{black!70}
      #2
    \end{minipage}
  \end{center}

  \vfill
  \begin{center}
    \small\color{black!60}Last updated: \today
  \end{center}
\end{titlepage}

% -----------------------------------------------------------------------------
% Table of contents -- kept on its own page
% -----------------------------------------------------------------------------
\clearpage
\setcounter{page}{2}
\thispagestyle{empty}
\hypersetup{linkcolor=black}
\setcounter{tocdepth}{2}
\tableofcontents
\hypersetup{linkcolor=brandprimary}
\clearpage
\pagenumbering{arabic}
}


\begin{document}
\sloppy
\makecover{Quiz Review: Theory}{Condensed theory from Lectures 1--3
(performance \& cost, instruction set architecture, pipelining): the
definitions, facts, comparisons and reasons most likely to be asked, followed
by a self-check question bank with answers.}

% ============================================================================
\section{Lecture 1: Computer Architecture, Performance and Cost}

\subsection{What computer architecture is}
\begin{definition}[Computer architect]
Designs and engineers the levels of a computer system to maximize
\textbf{performance} and \textbf{programmability} within the limits of
\textbf{technology} and \textbf{cost}. Tasks: instruction set design,
functional organization, logic design, implementation (IC design, packaging,
power, cooling).
\end{definition}
\begin{itemize}[leftmargin=*]
  \item The focus has shifted: arithmetic (1950s--60s) $\to$ ISA design
  (70s--80s) $\to$ CPU/memory/I/O/multiprocessors (90s) $\to$ multicore,
  on-chip networks, power (2000s) $\to$ low-power/mobile (2010s) $\to$ AI chips,
  quantum (2020s).
  \item \textbf{Levels of abstraction}: software (applications, compiler,
  assembler, OS, linker/loader, drivers) above the \textbf{ISA}, hardware
  (processor, memory, I/O, datapath \& control, logic, circuits, layout)
  below. The ISA is the hardware/software interface: one ISA, many
  implementations.
  \item \textbf{Von Neumann / stored-program computer} (EDSAC, 1949): the
  same memory holds instructions and data -- still the basic design today.
\end{itemize}

\subsection{Technology trends and the multicore shift}
\begin{itemize}[leftmargin=*]
  \item \textbf{Moore's Law}: transistors per cost-effective chip double
  every 18--24 months.
  \item Uniprocessor performance grew $\sim$25\%/yr (VAX era), $\sim$52\%/yr
  (RISC/x86, 1986--2002), then slowed sharply.
  \item \textbf{Power wall}: power is expensive, transistors are cheap.
  \textbf{ILP wall}: diminishing returns from more hardware for
  instruction-level parallelism. \textbf{Memory wall}: memory is slow
  ($\sim$200 cycles) relative to arithmetic. Together: the \textbf{brick
  wall}.
  \item Result: industry moved to \textbf{multicore} -- many simpler cores are
  more power-efficient than one complex core (``the processor is the new
  transistor''). The hard part is software: programming many cores.
\end{itemize}

\subsection{The design cycle}
Evaluate existing systems for bottlenecks (\textbf{performance}, using
\emph{benchmarks}) $\to$ simulate new designs (\textbf{technology}, using
\emph{workloads}) $\to$ implement the next generation (\textbf{cost},
\emph{implementation complexity}) $\to$ repeat.

\subsection{Measuring performance}
\begin{definition}[Latency vs.\ throughput]
\textbf{Latency} (execution/response time): time to finish one task.
\textbf{Throughput} (bandwidth): tasks completed per unit time. They often
conflict (Concorde has lower latency, the 747 higher throughput).
\end{definition}
\begin{itemize}[leftmargin=*]
  \item \textbf{CPU time} excludes I/O wait and other programs; it is user
  CPU time $+$ system CPU time.
  \item $\text{Performance} = 1/\text{ExTime}$; Speedup $=
  \text{ExTime}_{old}/\text{ExTime}_{new}$; ``X is $n\%$ faster than Y'':
  $\text{ExTime}_Y/\text{ExTime}_X = 1+n/100$.
  \item \textbf{CPU equation}: $T = I\times\CPI\times C$. Program, compiler
  and ISA affect $I$ and CPI; organization affects CPI and $C$; technology
  affects $C$.
  \item \textbf{Execution time of real programs is the only reliable
  measure}. MIPS, MFLOPS, clock rate, cache size alone are not.
\end{itemize}

\begin{theorem}[Amdahl's Law]
Speedup $=1/\big((1-F)+F/S\big)$: the improvement is limited by the fraction
of time $F$ the enhancement can be used; even with $S\to\infty$, speedup
$\le 1/(1-F)$.
\end{theorem}
\begin{note}
\textbf{Make the common case fast.} Every instruction is fetched but only
some access data $\Rightarrow$ optimize instruction fetch. Programs show
\textbf{locality} (90\% of time in 10\% of code; \emph{temporal}: reuse
soon, \emph{spatial}: neighbours soon) $\Rightarrow$ a memory
\textbf{hierarchy} (registers, cache, memory, disk) keeps frequent accesses
in small fast memories.
\end{note}

\subsection{MIPS and MFLOPS -- why they mislead}
\begin{itemize}[leftmargin=*]
  \item MIPS $=I/(T\cdot10^6)=f/(\CPI\cdot10^6)$ ignores instruction count:
  (1) depends on the ISA, (2) varies between programs on one machine, (3) can
  vary \emph{inversely} with performance. Valid only for the same program,
  ISA and compiler. \emph{Peak} MIPS uses only the lowest-CPI instructions.
  \item MFLOPS: better across different ISAs, but program-dependent (a
  compiler has $\approx0$ FP ops) and dependent on the FP operation mix;
  \emph{peak} MFLOPS is unrepresentative.
\end{itemize}

\subsection{Benchmarks and summarizing}
\begin{itemize}[leftmargin=*]
  \item Benchmark program levels, most to least accurate: \textbf{real
  applications, modified applications, kernels, toy benchmarks, synthetic
  benchmarks}.
  \item \textbf{SPEC} (System Performance Evaluation Cooperative): SPEC89
  (single SPECmark), SPEC92 (int/fp split), SPEC95, SPEC2000, SPEC CPU2006;
  SPECSFS and SPECWeb for servers.
  \item Arithmetic mean tracks total time but favours long programs; weighted
  mean depends on the chosen weights; \textbf{geometric mean} of normalized
  times is consistent whatever the reference machine, but does not track
  total time; harmonic mean averages rates.
  \item ``Benchmarks shape a field'': vendors may tune for benchmarks rather
  than real programs (``sales almost always wins'').
\end{itemize}

\subsection{Cost}
\begin{itemize}[leftmargin=*]
  \item \textbf{Cost} = what it takes to build; \textbf{price} = what it
  sells for. List price $=$ component cost $+$ direct cost (labour, scrap,
  warranty) $+$ gross margin (R\&D, marketing, profit) $+$ average discount.
  \item IC cost is driven by \textbf{yield}, \textbf{volume} and
  \textbf{feature size}. Larger dies: fewer dies per wafer \emph{and} lower
  yield, so die cost grows faster than die area.
  \item Learning curve: cost drops $\approx10\%$ each time volume doubles.
  High-end machines sell in low volume, so they need high margins to recoup
  R\&D.
\end{itemize}

% ============================================================================
\section{Lecture 2: Instruction Set Architecture}

\subsection{Basics}
\begin{itemize}[leftmargin=*]
  \item \textbf{Instruction set}: all instructions a processor can execute;
  the interface between software and hardware.
  \item \textbf{Execution cycle}: instruction fetch $\to$ decode $\to$
  operand fetch $\to$ execute $\to$ write back (store result).
  \item Instruction fields: \textbf{opcode}, \textbf{source operand
  reference(s)}, \textbf{result operand reference}.
  \item Design aspects: operation repertoire, addressing modes, data types,
  instruction format, registers.
  \item Instruction types: \textbf{data processing} (arithmetic, logic),
  \textbf{data storage} (load/store memory), \textbf{data movement}
  (register/I/O), \textbf{control} (test, branch).
\end{itemize}

\subsection{ISA styles (where the operands are)}
\begin{center}
\small
\begin{tabular}{@{}lp{10cm}@{}}
\toprule
Stack & operands on top of stack; dense code, simple HW; stack bottleneck, little parallelism \\
Accumulator & one implicit register; few instructions; heavy memory traffic, variable CPI \\
Memory-memory & all operands in memory; simple but memory-bound \\
Register-memory & one operand may be in memory (x86, 68000) \\
Load-store & ALU ops on registers only; \texttt{load}/\texttt{store} access memory (MIPS, PowerPC, Alpha) -- the modern choice \\
\bottomrule
\end{tabular}
\end{center}
Other styles: high-level-language architectures (Burroughs B5000) failed for
lack of compilers able to exploit them; RISC simplifies hardware and lets the
compiler build complex operations.

\subsection{Addressing modes}
\begin{definition}[Addressing mode]
How the processor refers to an operand (or a jump target). Classified by the
source of data (immediate, register, memory) and the address calculation
(direct, indirect, indexed).
\end{definition}
Immediate, register, direct, register indirect, displacement
(base+offset), indexed (base+index), memory indirect, PC-relative, scaled.
Measurements: \textbf{immediate and displacement are by far the most used};
most displacements fit in about 12--16 bits, so 16-bit fields suffice.

\subsection{Operands, encoding, registers}
\begin{itemize}[leftmargin=*]
  \item Operand types: characters (ASCII 8-bit, Unicode 16-bit), integers
  (2's complement, 16/32-bit), floating point (single 32-bit, double
  64-bit), BCD (business), fixed point (DSP), RGBA (graphics).
  \item Encoding: \textbf{variable length} (VAX, x86: dense code, hard
  decode), \textbf{fixed length} (MIPS, ARM, PowerPC: easy decode, bigger
  code), \textbf{hybrid} (MIPS16, Thumb: two fixed lengths).
  \item Registers are the fastest storage and flexible for compilers, but
  more registers need more encoding bits and a bigger register file can be
  slower.
\end{itemize}

\subsection{CISC vs.\ RISC}
\begin{center}
\small
\begin{tabular}{@{}p{6.9cm}p{6.9cm}@{}}
\toprule
CISC & RISC \\
\midrule
More work per instruction, compact code & Fewer, simpler instructions; more instructions executed \\
Born when memory was very expensive & Designed for pipelining, CPI $\le1$ \\
Many addressing modes (68000: 14; 68020: 25; 680x0 listing: 18; IA-32: 12) & Few modes (Alpha: 4; HP PA-RISC: 7) \\
Variable-length instructions (x86: 1--12 bytes) & Fixed 32-bit instructions \\
Memory operands in ALU instructions; backward compatibility & Load-store; delayed loads/branches; prefetch and speculation \\
\bottomrule
\end{tabular}
\end{center}

\subsection{MIPS}
\begin{itemize}[leftmargin=*]
  \item 32 general registers, plus PC, Hi, Lo; three 32-bit formats:
  \textbf{R} (op rs rt rd shamt funct), \textbf{I} (op rs rt imm16),
  \textbf{J} (op addr26); opcode always in the first 6 bits.
  \item \textbf{Design principles}: simplicity favours regularity; smaller
  is faster; make the common case fast; good design demands good compromises
  (three formats, all the same length).
  \item SPEC2006 mix: integer code is dominated by data transfer (35\%) and
  conditional branches (34\%); FP code by arithmetic (48\%) and data
  transfer (36\%).
  \item \textbf{Big-endian} (most significant byte at the lowest address).
  \item MIPS addressing: register, immediate, base+displacement
  (\texttt{lw}/\texttt{sw}), PC-relative (\texttt{beq}/\texttt{bne}, offset in
  words from PC$+4$), pseudo-direct (\texttt{j}/\texttt{jal}), register
  indirect (\texttt{jr}).
  \item No \texttt{MOV}, \texttt{NOT}, \texttt{SUBI}, \texttt{SLA}: use
  \texttt{add \dots,\$zero}, \texttt{nor \dots,\$zero}, \texttt{addi} with a
  negative constant, \texttt{sll}. 32-bit constants: \texttt{lui} $+$
  \texttt{ori}.
  \item \textbf{Atomic read-modify-write}: \texttt{ll}/\texttt{sc};
  \texttt{sc} fails (returns 0) if the word changed since the \texttt{ll}.
  \item \textbf{Procedure call}: save return address, save needed registers,
  execute, restore registers, restore return address, return.
  \texttt{jal} sets \texttt{\$ra}; \texttt{jr \$ra} returns. \texttt{\$s0--\$s7,
  \$sp, \$fp, \$gp, \$ra} are preserved across calls (callee saves);
  \texttt{\$t}, \texttt{\$a}, \texttt{\$v} are not (caller saves if needed). A
  leaf procedure need not save \texttt{\$ra}.
\end{itemize}

% ============================================================================
\section{Lecture 3: Pipelining}

\subsection{Concepts}
\begin{definition}[Pipelining]
Overlapping the execution of multiple instructions. Each step is a
\textbf{pipeline stage}; stages are connected linearly; the \textbf{machine
cycle} is set by the slowest stage.
\end{definition}
\begin{itemize}[leftmargin=*]
  \item Increases \textbf{throughput}; does \textbf{not} reduce (may slightly
  increase) the \textbf{latency} of one instruction ($k$ cycles).
  \item Ideal CPI $=1$: one instruction completes per cycle once the pipeline
  is full ($\CPI = (k+n-1)/n\to1$).
  \item Ideal speedup $=$ number of stages. Reduced by unbalanced stages,
  fill/drain time, latch/control overhead, and hazards.
  \item \textbf{Laundry lessons}: rate limited by the slowest stage; tasks
  overlap using different resources; dependencies force waiting.
\end{itemize}

\subsection{The 5-stage MIPS pipeline}
\begin{itemize}[leftmargin=*]
  \item IF (instruction memory) -- ID (decode, read registers) -- EX (ALU,
  address) -- MEM (data memory) -- WB (write register). Pipeline registers
  IF/ID, ID/EX, EX/MEM, MEM/WB carry values and control between stages.
  \item Register file: \textbf{write in the first half, read in the second
  half} of the cycle.
  \item \textbf{Why MIPS is easy to pipeline}: (1) all instructions the same
  length, (2) few formats with register fields in fixed positions (read
  registers while decoding), (3) memory operands only in \texttt{lw}/\texttt{sw}
  (EX computes the address, MEM accesses), (4) aligned operands (one access
  per transfer).
\end{itemize}

\subsection{Hazards}
\begin{definition}[Hazard]
A situation that prevents the next instruction from executing in its
designated cycle, forcing stalls and raising CPI above 1:
\textbf{structural} (resource conflict), \textbf{data} (dependence on a
previous result), \textbf{control} (branches and other PC changes).
\end{definition}
\begin{itemize}[leftmargin=*]
  \item $\CPI_{pipe} = 1 + $ stall cycles per instruction; speedup $=$
  depth$/(1+$stall CPI$)$ for balanced stages.
  \item \textbf{Structural}: e.g.\ one memory port for instructions and
  data. Fix by duplicating resources (separate I/D memories) or stalling;
  sometimes accepted when rare and costly to remove.
  \item \textbf{Data hazard classes} ($I$ before $J$): \textbf{RAW} (true
  dependence; $J$ reads too early), \textbf{WAW} and \textbf{WAR} (name
  dependences; writes in the wrong order / $J$ writes before $I$ reads);
  RAR is not a hazard. Only RAW occurs in the simple in-order 5-stage
  pipeline.
  \item \textbf{Forwarding (bypassing)}: send a result from the EX/MEM or
  MEM/WB pipeline register straight to the ALU input, instead of waiting for
  write-back.
  \item \textbf{Load-use hazard}: a load's data is ready only after MEM, so
  the next instruction still stalls one cycle even with forwarding.
  \item \textbf{Hardware interlock}: detects the hazard and stalls.
  \textbf{Compiler (static) scheduling}: reorders code so a load is not
  immediately followed by its use. Static $=$ compile time, dynamic $=$ run
  time in hardware.
\end{itemize}

\subsection{Control hazards and branches}
\begin{itemize}[leftmargin=*]
  \item Branch penalty $=$ (stage where the branch is resolved) $-1$: 3 cycles
  if resolved in MEM. MIPS moves the zero test and target adder into ID
  $\Rightarrow$ 1 cycle.
  \item \textbf{Four alternatives}: (1) stall until resolved; (2)
  \textbf{predict not taken} -- keep fetching, squash if taken (late state
  update makes this safe; $\approx47\%$ of MIPS branches not taken); (3)
  \textbf{predict taken} ($\approx53\%$ taken) -- no gain in MIPS because the
  target is not known earlier; (4) \textbf{delayed branch} -- the
  instruction in the delay slot always executes (MIPS uses one slot).
  \item Delay-slot filling: from \textbf{before} the branch (best, always
  useful, branch must not depend on it); from the \textbf{target} (useful if
  taken, may need copying, must be safe if not taken); from the
  \textbf{fall-through} (useful if not taken, must be safe if taken).
  \textbf{Cancelling (nullifying) branches} make the last two easier.
  \item Compilers fill $\approx60\%$ of slots, $\approx80\%$ of them usefully
  ($\approx50\%$ overall). Delayed branches lost popularity: deeper and
  multiple-issue pipelines need more slots, and dynamic prediction became
  cheap.
\end{itemize}

% ============================================================================
\section{Self-Check Question Bank}

Answers are in the next section.

\begin{enumerate}[label=\textbf{Q\arabic*.},leftmargin=*]
  \item What is the role of a computer architect?
  \item Why is the ISA called the hardware/software interface?
  \item Name the three ``walls'' and the conclusion drawn from them.
  \item Latency vs.\ throughput: which one does pipelining improve?
  \item Which of $I$, CPI, $C$ does (a) the compiler, (b) the organization,
  (c) the technology affect?
  \item State Amdahl's Law and its main consequence.
  \item Give three reasons why MIPS is a poor performance metric. When is it
  valid?
  \item Order the five kinds of benchmark programs from most to least
  accurate.
  \item Why use the geometric mean of normalized times? What is its weakness?
  \item Why does die cost grow faster than die area?
  \item Cost vs.\ price: what is between them?
  \item List the five steps of the instruction execution cycle.
  \item Compare stack and accumulator ISAs (one advantage and one drawback
  each).
  \item Which addressing modes are used most, and what does that imply for
  instruction design?
  \item Fixed vs.\ variable-length encoding: trade-off?
  \item Give four characteristics of RISC.
  \item State the four MIPS design principles.
  \item How do you perform \texttt{NOT}, \texttt{MOV}, and load a 32-bit
  constant in MIPS?
  \item What do \texttt{ll}/\texttt{sc} provide?
  \item Which registers must a callee preserve? Must a leaf procedure save
  \texttt{\$ra}?
  \item Why is the ideal CPI of a pipeline 1 even though each instruction
  takes 5 cycles?
  \item What limits pipeline speedup in practice?
  \item Give four reasons why MIPS is easy to pipeline.
  \item Define the three kinds of hazard with an example of each.
  \item Why can't WAR and WAW hazards occur in the 5-stage MIPS pipeline?
  \item What is forwarding, and which hazard can it not remove?
  \item How does MIPS reduce the branch penalty from 3 to 1 cycle?
  \item Compare predict-not-taken and delayed branch.
  \item What are the three ways to fill a delay slot and their conditions?
  \item Why have delayed branches fallen out of favour?
\end{enumerate}

% ============================================================================
\section{Answers}

\begin{enumerate}[label=\textbf{A\arabic*.},leftmargin=*]
  \item Design the levels of a computer system to maximize performance and
  programmability within technology and cost limits.
  \item It is the only layer both sides agree on: software is written to it
  and every hardware implementation must execute it, so one ISA supports
  many implementations.
  \item Power wall, ILP wall, memory wall (together the brick wall)
  $\Rightarrow$ single-core performance growth slowed, so chips went
  multicore.
  \item Pipelining improves throughput, not latency.
  \item (a) $I$ and CPI; (b) CPI and $C$; (c) $C$.
  \item Speedup $=1/((1-F)+F/S)$; the part not improved limits the overall
  gain (at most $1/(1-F)$) $\Rightarrow$ make the common case fast.
  \item ISA-dependent; varies between programs; can vary inversely with
  performance (ignores instruction count). Valid only with the same program,
  ISA and compiler.
  \item Real applications, modified applications, kernels, toy benchmarks,
  synthetic benchmarks.
  \item Its ranking does not depend on the reference machine. It does not
  reflect total execution time.
  \item A larger die means fewer dies per wafer and a lower yield (more
  likely to contain a defect), so cost rises faster than area.
  \item Direct costs, gross margin and the average discount are added to the
  component cost to give the list price.
  \item Fetch, decode, operand fetch, execute, write back.
  \item Stack: dense code / stack bottleneck. Accumulator: few instructions,
  easy compiler / heavy memory traffic, variable CPI.
  \item Immediate and displacement; so a good ISA supports them efficiently
  with $\approx16$-bit fields.
  \item Fixed: simple, fast decode (good for pipelining) but larger code;
  variable: compact code but complex decode.
  \item Simple instructions, fixed-length encoding, load-store, few
  addressing modes, designed for pipelining with CPI $\approx1$ (also
  delayed branches/loads).
  \item Simplicity favours regularity; smaller is faster; make the common
  case fast; good design demands good compromises.
  \item \texttt{nor \$d,\$s,\$zero}; \texttt{add \$d,\$s,\$zero};
  \texttt{lui} upper half then \texttt{ori} lower half.
  \item An atomic read-modify-write: \texttt{sc} succeeds only if the word was
  not modified since the \texttt{ll}.
  \item \texttt{\$s0--\$s7}, \texttt{\$sp}, \texttt{\$fp}, \texttt{\$gp},
  \texttt{\$ra}. No -- a leaf never executes \texttt{jal}, so \texttt{\$ra} is
  untouched.
  \item Once full, one instruction finishes every cycle; the $k-1$ fill cycles
  are paid once, so $\CPI=(k+n-1)/n\to1$.
  \item Unbalanced stages, pipeline-register overhead, fill/drain time, and
  hazards (stalls).
  \item Same instruction length; few formats with fixed register fields;
  memory operands only in loads/stores; aligned operands.
  \item Structural (one memory port for instruction and data accesses), data
  (\texttt{add r1}; \texttt{sub} using \texttt{r1}), control (a taken
  \texttt{beq}).
  \item Reads happen early (ID) and writes late (WB), and instructions stay in
  order, so a later instruction can neither write before an earlier one reads
  nor before an earlier one writes.
  \item Routing a result from a pipeline register directly to the ALU input
  of a later instruction. It cannot remove the load-use stall.
  \item It tests the register in ID and computes the target with an extra
  adder in ID, so the branch is resolved in stage 2.
  \item Predict-not-taken fetches sequentially and squashes on a taken branch
  (penalty only when taken); delayed branch always executes the slot
  instruction (cost only when the slot holds no useful work).
  \item From before the branch (independent of the branch); from the target
  (safe if not taken, may duplicate code); from the fall-through (safe if
  taken).
  \item Deeper and multiple-issue pipelines need more than one slot, and
  dynamic branch prediction became cheap.
\end{enumerate}

\end{document}
